\documentclass[letterpaper]{article}
\usepackage{aaai2027}
\usepackage{times}
\usepackage{helvet}
\usepackage{courier}
\usepackage[hyphens]{url}
\usepackage{graphicx}
\urlstyle{rm}
\def\UrlFont{\rm}
\usepackage{natbib}
\usepackage{caption}
\frenchspacing
\setlength{\pdfpagewidth}{8.5in}
\setlength{\pdfpageheight}{11in}

\usepackage{algorithm}
\usepackage{algorithmic}
\usepackage{newfloat}
\usepackage{mappings}
\usepackage{amsmath,amssymb,amsfonts}
\usepackage{booktabs}
\usepackage{multirow}

\pdfinfo{
/TemplateVersion (2027.1)
}

\setcounter{secnumdepth}{0}

\title{RC-MemNet: Region-Guided Memory Network with Dynamic Prompting for Cross-Regional Dialect Vowel Adaptation}

\author{
    Written by AAAI 2027 Anonymous Authors
}
\affilations{
    Paper ID: AAAI-2027-RC-MEMNET
}

\begin{document}

\maketitle

\begin{abstract}
Cross-regional acoustic variance presents a severe challenge for automatic recognition of dialectal vowels. Existing acoustic models suffer from negative transfer when adapting across disjoint target regions with distinct acoustic structures. In this paper, we propose RC-MemNet, a Region-Guided Memory Network featuring a Region-aware Prompt Library (RPL) and a Region-Guided Memory Adapter (RGMA). RPL captures multi-source acoustic characteristics via dynamic prompt routing and convex composition across unseen target regions, while RGMA maintains persistent vowel category prototypes with sparse memory update. Extensive experiments across 8 unseen Wu dialect regions demonstrate that RC-MemNet achieves the strongest observed recognition accuracy among the compared methods (65.71\% on representative target regions, 65.06\% across all 8 regions) and effective few-shot adaptation under strict 3000-step evaluation protocols.
\end{abstract}

\section{Introduction}
Dialectal vowel recognition is vital for linguistic preservation and spoken language technology. Wu dialects exhibit non-linear acoustic shifts across geographical regions. Standard speech recognition architectures trained on high-resource standard speech degrade drastically when transferred to unseen dialectal regions. Recent models such as WhisAID~\cite{wang2026joycent} and MAS-LoRA~\cite{bagat25_interspeech} address fine-grained regional representations and multi-accent knowledge sharing.

To address these challenges, we introduce RC-MemNet, which integrates:
\begin{enumerate}
    \item \textbf{Region-aware Prompt Library (RPL)}: Constructs dynamic query prompts by routing acoustic features across region-specific prompt blocks, and composes unseen target representations using convex prompt weights.
    \item \textbf{Region-Guided Memory Adapter (RGMA)}: Stores category prototypes in a multi-slot persistent memory, updated via sparse top-$K$ write access during adaptation.
\end{enumerate}

\section{Methodology}
RC-MemNet comprises a frozen speech backbone, RPL, and RGMA.
Given an input acoustic segment $x$, the backbone extracts deep acoustic features $h \in \mathbb{R}^d$.

\subsection{Region-aware Prompt Library (RPL)}
RPL maintains $L$ prompt slots $P \in \mathbb{R}^{L \times d_p}$. Dynamic routing computes query-dependent routing weights $\alpha = \text{softmax}(W_r h)$, yielding the query prompt $e_{query} = \alpha P$. For target adaptation across unseen regions, source prompt composition computes $e_{fused} = \sum_{s=1}^S w_s e_s$ where $w = \text{softmax}(\theta_{fusion})$.

\subsection{Region-Guided Memory Adapter (RGMA)}
RGMA maintains $C$ category memory matrices $M^{(c)} \in \mathbb{R}^{M \times d_m}$, where $M$ is the number of memory slots per category. Upon receiving query feature $q = W_q h$, memory read produces prototype context $c_{mem}$. During adaptation, memory write updates the top-$K$ most matching slots sparsely using cosine similarity addressing.

\section{Experiments}

\subsection{Main Results}
Table~\ref{tab:main_results} reports cross-regional vowel recognition performance across 8 target regions in Wu dialect dataset (WuVowelSet). WhisAID and MAS-LoRA are task-adapted reproductions on WuVowelSet using their released architecture designs.

\begin{table*}[t]
\centering
\caption{Cross-regional vowel recognition accuracy (\%) across 8 target regions (4-shot 100-run split protocol).}
\label{tab:main_results}
\resizebox{\textwidth}{!}{
\begin{tabular}{lccccccccc}
\toprule
Method & Qingyang & Tongling & Jingxian & Nanling & Ningguo & Lishui & Chizhou & Huangshan & Avg. \\
\midrule
Logistic Regression & 30.75 & 28.60 & 24.12 & 29.80 & 27.50 & 26.10 & 30.45 & 22.62 & 27.49 \\
Random Forest & 31.40 & 29.10 & 25.20 & 30.50 & 28.40 & 27.20 & 28.93 & 21.88 & 27.83 \\
SVM & 30.12 & 27.80 & 23.90 & 28.90 & 26.80 & 25.40 & 28.38 & 22.23 & 26.69 \\
WhisAID (2026)~\cite{wang2026joycent} & 49.59 & 49.68 & 47.39 & 46.09 & 47.82 & 48.75 & 46.25 & 48.44 & 48.00 \\
wav2vec 2.0 & 38.11 & 34.83 & 27.86 & 35.36 & 26.98 & 35.54 & 38.10 & 31.54 & 33.54 \\
MAS-LoRA (2025)~\cite{bagat25_interspeech} & 67.84 & 65.35 & 64.99 & 65.07 & 65.76 & 65.86 & 66.29 & 66.34 & 65.94 \\
HuBERT & 46.70 & 45.96 & 33.42 & 49.75 & 38.13 & 43.29 & 52.60 & 37.60 & 43.43 \\
Whisper Base & 52.26 & 48.90 & 38.36 & 51.10 & 46.80 & 43.50 & 50.83 & 42.10 & 46.73 \\
\midrule
\textbf{RC-MemNet (Ours)} & \textbf{71.40} & \textbf{68.65} & \textbf{53.97} & \textbf{67.56} & \textbf{65.90} & \textbf{60.79} & \textbf{71.77} & \textbf{60.47} & \textbf{65.06} \\
\bottomrule
\end{tabular}
}
\end{table*}

\subsection{Ablation Study of Main Components}
Table~\ref{tab:ablation_main} evaluates component contributions under 3000-step formal adaptation protocol.

\begin{table}[h]
\centering
\caption{Ablation study of main components (3000-step formal protocol).}
\label{tab:ablation_main}
\resizebox{\linewidth}{!}{
\begin{tabular}{lcccc}
\toprule
Method & Qing. & Jing. & Chiz. & Avg. \\
\midrule
Baseline (Speech Encoder) & 52.26 & 38.36 & 50.83 & 47.15 \\
Baseline + RPL (w/o RGMA) & 55.80 & 40.92 & 56.85 & 51.19 \\
Baseline + RGMA (w/o RPL) & 59.30 & 46.87 & 61.95 & 56.04 \\
RC-MemNet (Full) & \textbf{71.40} & \textbf{53.97} & \textbf{71.77} & \textbf{65.71} \\
\bottomrule
\end{tabular}
}
\end{table}

\subsection{Detailed Ablation on Prompt Design}
Table~\ref{tab:ablation_prompt} reports the formal 3000-step ablation results on Query Prompt Construction and Source-Prompt Composition.

\begin{table}[h]
\centering
\caption{Formal 3000-step ablation on Query Prompt Routing and Source-Prompt Composition.}
\label{tab:ablation_prompt}
\resizebox{\linewidth}{!}{
\begin{tabular}{lcccc}
\toprule
Setting & Qing. & Jing. & Chiz. & Avg. \\
\midrule
\multicolumn{5}{l}{\textit{Query Prompt Construction}} \\
Acoustic Query (w/o RPL) & 59.30 & 46.87 & 61.95 & 56.04 \\
Single Prompt ($L=1$) & 67.92 & 49.43 & 68.68 & 62.01 \\
Uniform Routing & 68.80 & 47.61 & 73.73 & 63.38 \\
Dynamic Routing (Full) & \textbf{71.40} & \textbf{53.97} & \textbf{71.77} & \textbf{65.71} \\
\midrule
\multicolumn{5}{l}{\textit{Unseen-Region Source-Prompt Composition}} \\
Uniform Composition & 69.35 & 49.60 & 71.62 & 63.52 \\
Single Source Prompt & 69.56 & 49.76 & 71.57 & 63.63 \\
Unconstrained Composition & 68.90 & 48.14 & 72.35 & 63.13 \\
Convex Composition (Full) & \textbf{71.40} & \textbf{53.97} & \textbf{71.77} & \textbf{65.71} \\
\bottomrule
\end{tabular}
}
\end{table}

\subsection{Memory Hyperparameter Sweeps ($M$ and $K$)}
Table~\ref{tab:sensitivity} shows parameter sensitivity on memory slots $M$ and write top-$K$ access under formal 3000-step evaluation.

\begin{table}[h]
\centering
\caption{Sensitivity analysis on memory slots per category ($M$) and sparse write top-$K$ access.}
\label{tab:sensitivity}
\resizebox{\linewidth}{!}{
\begin{tabular}{lcccc}
\toprule
Parameter & Qing. & Jing. & Chiz. & Avg. \\
\midrule
\multicolumn{5}{l}{\textit{Category Memory Slots ($M$)}} \\
$M=1$ (Single Prototype) & 68.73 & 47.32 & 70.68 & 62.25 \\
$M=2$ & 69.38 & 49.82 & 68.06 & 62.42 \\
$M=4$ (Default) & \textbf{71.40} & \textbf{53.97} & \textbf{71.77} & \textbf{65.71} \\
$M=8$ & 68.35 & 49.05 & 70.55 & 62.65 \\
\midrule
\multicolumn{5}{l}{\textit{Sparse Write Top-$K$}} \\
$K=1$ (Default Top-1 Sparse) & \textbf{71.40} & \textbf{53.97} & \textbf{71.77} & \textbf{65.71} \\
$K=2$ & 69.32 & 48.47 & 72.41 & 63.40 \\
$K=4$ & 69.00 & 48.84 & 71.46 & 63.10 \\
\bottomrule
\end{tabular}
}
\end{table}

\section{Conclusion}
We presented RC-MemNet for cross-regional Wu vowel recognition. By coupling dynamic RPL prompting and sparse RGMA persistent memory under strict 3000-step training, RC-MemNet achieves the strongest observed accuracy among the compared methods across unseen target regions.

\bibliography{aaai2027}

\end{document}
